% Audited unnumbered auxiliary.  Ordinary equivalence is a separate endpoint problem.
\paragraph{The rational Chow group of the punctured tetrahedral germ
(audited auxiliary calculation).}
Let $k_0=\overline{\mathbf F}_p$, let $k/k_0$ be algebraically closed,
and take the exact ordered tetrahedral construction TG.1--TG.5 over
$k_0$, followed by base change to $k$ and strict henselization at its
vertex.  Write $R$ for that strictly henselian normal local ring,
$\pi:Z\to\operatorname{Spec}R$ for the regular projective resolution,
$E=\bigcup_{i=0}^3E_i$ for its reduced exceptional divisor, and
$X=Z\setminus E$.  Then, after choosing the orientation and splitting
of TS.2, there is an isomorphism
\begin{equation}
                 \mathrm{CH}^{2}(X)_{\mathbf Q}
                       \simeq K_2(k)_{\mathbf Q}.
                 \tag{TC.1}
\end{equation}
This statement concerns rational equivalence.  It neither identifies
the ordinary algebraically trivial subgroup nor computes
$B^2(X)_{\mathbf Q}$.

\emph{Step 1: the weight-two group on the regular resolution.}
By the point-parameter isomorphism FF.4, using Drinfeld's
Theorem~3.7 and Remark~(ii)\citep[author p.~8]{frDrinfeld2004Tate},
restriction induces
\[
                     K_0(Z)_{\mathbf Q}
                            \xrightarrow{\sim}K_0(E)_{\mathbf Q}.
\]
It commutes with Adams operations.  The rational Chern character
identifies $\mathrm{CH}^{2}(Z)_{\mathbf Q}$ with the weight-two
part of $K_0(Z)_{\mathbf Q}$.  Here is the finite-type justification
in the present strict-local setting.  The original resolution is
defined over a finite-type $k$-scheme; its base changes over pointed
etale neighbourhoods of the vertex are regular after shrinking
around the proper exceptional fibre.  Since $k$ is perfect, these
regular finite-type models are smooth.  Every finite collection of
perfect complexes, cycles, rational equivalences, Chern classes and
relations descends to one such neighbourhood.  The rational Chern
character isomorphism on those smooth models
\citep[Example~15.2.16, pp.~294--295]{frC4Fulton1998}
therefore passes to the filtered strict-local limit, as in C.3ad.

The whole-diagram comparison of TS, with its compatible Adams
decomposition, now gives
\begin{equation}
 \mathrm{CH}^{2}(Z)_{\mathbf Q}
       \simeq \mathbf Q^4\oplus K_2(k)_{\mathbf Q}.
                                                        \tag{TC.2}
\end{equation}
The first four coordinates are restriction to each $E_i$ followed
by the degree of its zero-cycle Chern character.  The last coordinate
is the chosen $\kappa K_2(k)$ coefficient from TS.4.  On a field,
$K_2(k)$ has Adams weight two
\citep[Chapter~IV, Example~5.9.1, p.~51]{D_Weibel2013KbookIV}.
There is no support or Thom shift here: the operations on the closed
restriction diagram commute with its homotopy-fibre boundaries.
The exact-category construction and its scope were spelled out in
C.3bc, Step~2; this does not apply a regular-ambient theorem directly
to a singular scheme.

\emph{Step 2: exceptional curves have zero coefficient in $K_2(k)$.}
The cycle map from the disjoint union of the four components gives
a surjection
\[
             \bigoplus_i\mathrm{CH}_1(E_i)
                         \longrightarrow\mathrm{CH}_1(E).
\]
Each $E_i$ is a blowup of $\mathbf P^2$ in $15,10,5,0$ specified
$k_0$-points, respectively.  Its divisor group modulo rational
equivalence is freely generated by a pulled-back line and the
point-exceptional curves.  Every one of these generators is defined
over $k_0$.

Their pushes into $Z$, and subsequent restrictions to $E$ in
$K$-theory, are base changes of the same perfect proper-pushforward
and pullback operations on the regular finite-type construction
over $k_0$.  Indeed $Z$ is obtained from that construction by flat
base change and strict localization, so flat base change for these
perfect closed immersions applies.  The corresponding rational
class on $E_0$ has no $K_2$ coefficient, because
$K_2(k_0)_{\mathbf Q}=0$.  Naturality of TS.1 under this base change
therefore proves that every exceptional curve has zero last
coordinate in (TC.2).  Taking the weight-two Chern component gives
precisely the ordinary Chow pushforward of its divisor class.
This argument uses the actual $k_0$-defined generators; it does not
assume that $K(Z)$ is a tensor base change of $K(Z_{k_0})$.

\emph{Step 3: exceptional curves span all four degree coordinates.}
TG.5 supplies an effective exceptional divisor
$D=\sum_i a_iE_i$, with all $a_i>0$, such that
$\mathcal O_Z(-D)|_E$ is ample.  Choose a relatively ample line
bundle $\mathcal H$ for $Z/\operatorname{Spec}R$, and let
$H=c_1(\mathcal H)$.  It can be chosen from the construction over
$k_0$.  Set
\[
       M_{ij}=\deg(E_iE_jH),\qquad A=\operatorname{diag}(a_i).
\]
All intersections are proper numerical intersections supported on
the projective exceptional fibre.  The matrix $M$ is symmetric.
For $i\ne j$, $M_{ij}\geq0$, since $E_i\cap E_j$ is an effective
curve and $\mathcal H$ restricts amply to it.  Moreover
\[
       (Ma)_i=\deg(E_iDH)<0
\]
because both $-D|_{E_i}$ and $H|_{E_i}$ are ample on the projective
surface $E_i$.  For $B=AMA$ and every real vector $x$, the exact
identity
\[
 x^{\mathsf T}Bx
   =\sum_i\biggl(\sum_jB_{ij}\biggr)x_i^2
       -\sum_{i<j}B_{ij}(x_i-x_j)^2
\]
is strictly negative unless $x=0$: every row sum is negative and
every off-diagonal entry is nonnegative.  Thus $M$ is invertible.

For each $i$ take the exceptional curve class
$\gamma_i=c_1(\mathcal H|_{E_i})\cap[E_i]$ and push it to $Z$.
Its four restriction-degree coordinates are exactly the $i$th
column of $M$, by refined pullback and the Cartier intersection
formula.  Hence the exceptional image in (TC.2) is all of
$\mathbf Q^4\oplus0$, by Step~2.

\emph{Step 4: localization and the unresolved ordinary quotient.}
Chow localization gives the exact sequence
\[
       \mathrm{CH}_1(E)_{\mathbf Q}
           \longrightarrow\mathrm{CH}^{2}(Z)_{\mathbf Q}
           \longrightarrow\mathrm{CH}^{2}(X)_{\mathbf Q}
           \longrightarrow0.
\]
Steps~2--3 and (TC.2) prove (TC.1).  The degree convention is the
codimension convention inherited from the regular threefold; no
identification of residual transcendence degrees with Krull
dimension of the punctured topological space is needed.

For a fixed smooth proper curve $P/k$, TS gives a compatible
$\kappa K_2(P)$ coefficient on $E\times P$, but its lift to
$Z\times P$ is still governed by the boundary in FF.1--FF.3.
The proper-curve symbol lemma in the retained top-quotient research
shows that bare $K_2(k)$ coefficients can vary along actual proper
curves on $E$; it does not supply the missing lift on $Z$.
Accordingly (TC.1) is an exact reduction of this example's ordinary
quotient problem, not a proof that its displayed group survives
ordinary algebraic equivalence or vanishes modulo that equivalence.
The subsequent endpoint theorem EL supplies the additional
actual-family argument and settles that separate question for
this specified example.
